% !PRINTED 626
\weberpage{626}

Der Werth $u=0$ oder $v=0$ tritt nur dann ein, wenn
$\alpha=0$ ist. Dann haben wir in der That in (13) die wohlbekannte
 metacyklische Gleichung $x^5+\beta=0$; diesen Fall also
lassen wir jetzt bei Seite. Die Gleichung (20) kann nur dann
zwei entgegengesetzte Wurzeln haben, wenn $\Delta=0$ ist, wenn
also die Gleichung (13) gleiche Wurzeln hat. Auch dies ist
auszuschliessen, da wir (13) als irreducibel vorausgesetzt haben.

Es bleibt also nur die Frage \"ubrig, ob (20) gleiche Wurzeln
haben kann. Es m\"usste dann mit (20) zugleich die abgeleitete
Gleichung
\[
6u^5-20\alpha u^3+30\alpha^2u-\sqrt{\Delta}=0
\]
befriedigt sein, und wenn wir aus dieser und aus (20) $\sqrt{\Delta}$
eliminiren,
\[
5u^6-15\alpha u^4+15\alpha^2u^2-5\alpha^3=5(u^2-\alpha)^3=0.
\]
Es m\"usste also $v=\alpha$ sein und also nach (22) $\beta=0$.
Dann w\"are aber die Gleichung 5$^{\text{ten}}$ Grades wieder reducibel, da
sie den Factor $x$ hat.

Will man also eine gegebene Gleichung von der Form (13)
in Bezug auf ihre Aufl\"osbarkeit pr\"ufen, so hat man zuerst die
Irreducibilit\"at festzustellen, und dann zu untersuchen, ob (21)
oder (22) eine rationale Wurzel hat.

Sind $\alpha$ und $\beta$ ganze Zahlen, so muss auch ein rationales $v$
eine ganze Zahl sein, die unter den Factoren von $25\alpha^6$ zu
suchen ist.

Ist z.\,B. $\alpha=5$, $\beta=5t$, so ist, wenn $t$ eine durch 5 nicht
theilbare ganze Zahl ist, $x^5+\alpha x+\beta$ nach \S.\,178, 3. irreducibel
. F\"ur $v$ hat man Potenzen von 5, mit positivem oder
negativem Vorzeichen einzusetzen, aber f\"ur keine solche Zahl
kann (22) erf\"ullt sein, weil auf der linken Seite die Potenz von
5 nicht hoch genug wird.

\gesperrt{Es ist also keine Gleichung von der Form}
$x^5+5x+5t=0$ \gesperrt{metacyklisch}.

Will man metacyklische Gleichungen ermitteln, so setze man
in (22)
\[
\beta=\alpha\mu,\quad v=\alpha\lambda,
\]
wodurch man erh\"alt:
\weq{23}{\begin{aligned}
\alpha&=\dfrac{5^5\mu^4\lambda}{(\lambda-1)^4(\lambda^2+6\lambda+25)}\\
\beta&=\dfrac{5^5\mu^5\lambda}{(\lambda-1)^4(\lambda^2+6\lambda+25)},
\end{aligned}}
%CONT>>
